

\begin{frame}{Appendix --- Inference: Few Clusters and the Exposure-Robust Bound}
\hypertarget{app:inference}{}
\footnotesize
The point estimate is invariant to how standard errors are computed. What moves is precision, and with it the weak-IV verdict. We cluster at the state level, where the leave-own-state shift is constant.
\smallskip
\begin{itemize}\itemsep3pt
  \item \textbf{Why state.} The leave-own-state-out shifter is identical within a state, so state clustering is the conservative default. Residuals also correlate across units with similar shares wherever they are located \citep{adao2019}; the exposure-robust SE below addresses that.
  \item \textbf{Few clusters $\Rightarrow$ weak-IV-robust AR.} State clustering leaves only \textbf{\NClusters{} clusters}, where asymptotic cluster-robust SEs can under-cover. Headline inference is therefore carried by an Anderson--Rubin \textbf{wild-cluster restricted bootstrap} (\ns{\citealp{andersonrubin1949,roodman2019}}; $B=9999$, Rademacher weights over \ARNClusters{} states). It confirms the consolidation (margin $p=\ARMarginP$, runner-up share $p=\ARRunnerUpP$, both excluding zero), leaves the mayoral blank rate tentative (the AR set straddles zero, $p=\ARBlankP$), and leaves the null outcomes (turnout, winner-majority, female share) null.
  \item \textbf{Exposure-robust outer bound.} The BHJ/AKM standard error \ns{\citep{borusyak2022,adao2019}}, built for the few-shock shift-share structure and re-clustered at the subject (shock) level, inflates the headline margin SE by only $\ERMarginRatio\times$, so the consolidation still clears. The SE barely moves because the estimate does not hang on one dominant shock. \hfill\hyperlink{app:exposure}{\beamergotobutton{weak-IV \& exposure-robust appendix}}
\end{itemize}

\centerline{\hyperlink{src:identbrief}{\beamerreturnbutton{Back}}}
\end{frame}
